% Einheit 3 von 4 – Antiproportionale Zuordnungen (bank/zuordnungen/e3.jsonl, zusammenbau v0.1)
%% TODO Zweigzeile Teil 1: was hier gelernt wird (Satz in Schülersprache aus dem Einheitstitel) – Einheit 3
\einheitenkopf[e3]{Einheit 3 von 4 · Antiproportionale Zuordnungen}
\zweigzeile{ab Kl. 7 · P10}
\setcounter{aufgabe}{19}

%% TODO Titel: Ich-kann-Satz fehlt in der Bank (Platzhalter: Kettenname) – Nr. 20
%% TODO Anweisung: ein Satz über der ersten Teilaufgabe fehlt in der Bank – Nr. 20
\begin{aufgabe}{Antiproportional}
\begin{teile}
\teil Je mehr Maler an einer Wand arbeiten, desto … Zeit brauchen sie. \\ \kreuz{mehr} \\ \kreuz{weniger}
\end{teile}
%% TODO Darstellung für schwach fehlt in der Bank (zuordnungen-e3-k1-s1-v1; Abschnitt „Für schwache Schüler“)
\swz{2 Bagger brauchen für eine Baugrube 10 Tage. Wie lange brauchen 4 Bagger?}{}
%% TODO Darstellung für schwach fehlt in der Bank (zuordnungen-e3-k1-s1-v2; Abschnitt „Für schwache Schüler“)
\swz{Mit 40 km/h dauert eine Fahrt 30 min. Wie lange dauert sie mit 80 km/h?}{}
%% TODO Darstellung für schwach fehlt in der Bank (zuordnungen-e3-k1-s1-v3; Abschnitt „Für schwache Schüler“)
\swz{Wenn sich 4 Kinder eine Tüte Bonbons teilen, bekommt jedes 10 Stück. Wie viele bekommt jedes bei 8 Kindern?}{}
%% TODO Darstellung für schwach fehlt in der Bank (zuordnungen-e3-k1-s1-v4; Abschnitt „Für schwache Schüler“)
\swz{6 Maler brauchen für ein Haus 4 Tage. Wie lange brauchen 3 Maler?}{}
%% TODO Darstellung für schwach fehlt in der Bank (zuordnungen-e3-k1-s1-v5; Abschnitt „Für schwache Schüler“)
\swz{Packt man 20 Bücher in jeden Karton, braucht man 6 Kartons. Wie viele Kartons braucht man bei 40 Büchern je Karton?}{}
\swfrage{Was bleibt gleich, was ändert sich?}
%% TODO Darstellung für schwach fehlt in der Bank (zuordnungen-e3-k1-s2-v1; Abschnitt „Für schwache Schüler“)
\swz{2 Traktoren pflügen ein Feld in 9 Stunden. Wie lange brauchen 6 Traktoren?}{}
%% TODO Darstellung für schwach fehlt in der Bank (zuordnungen-e3-k1-s3-v1; Abschnitt „Für schwache Schüler“)
\swz{3 Maschinen brauchen für einen Auftrag 8 Stunden. Wie lange brauchen 4 Maschinen?}{}
%% TODO Darstellung für schwach fehlt in der Bank (zuordnungen-e3-k1-s4-v1; Abschnitt „Für schwache Schüler“)
\swz{Jemand behauptet: 6 Arbeiter brauchen 10 Tage, 4 Arbeiter brauchen 15 Tage. Prüfe mit dem Produkt, ob das zusammenpasst.}{}
%% TODO Darstellung für schwach fehlt in der Bank (zuordnungen-e3-k1-s5-v1; Abschnitt „Für schwache Schüler“)
\swz[3]{Für einen Umzug haben 3 Freunde 4 Stunden eingeplant. Kurzfristig helfen 3 weitere Freunde mit. Alle arbeiten gleich schnell. Wie lange dauert der Umzug jetzt?}{}
\swz[3]{Ein Futtervorrat reicht für 6 Ziegen 10 Tage. Wie viele Tage reicht er für 5 Ziegen? \hfill (P10 2014 GYM)}{}
\end{aufgabe}

%% TODO Titel: Ich-kann-Satz fehlt in der Bank (Platzhalter: Kettenname) – Nr. 21
%% TODO Anweisung: ein Satz über der ersten Teilaufgabe fehlt in der Bank – Nr. 21
\begin{aufgabe}{Graph als fallende Kurve, Werte ablesen}
\swz{Lies ab: Wie lange brauchen 3 Helfer? Wie viele Helfer braucht man, damit es 2 Stunden dauert?}{\begin{ksys}[xmin=0,xmax=7,ymin=0,ymax=13,ablesen,xlabel=Anzahl Helfer,ylabel=Zeit in h]\funktionab{12/\x}{}{1}{6}\end{ksys}}
\end{aufgabe}

%% TODO Titel: Ich-kann-Satz fehlt in der Bank (Platzhalter: Kettenname) – Nr. 22
%% TODO Anweisung: ein Satz über der ersten Teilaufgabe fehlt in der Bank – Nr. 22
\begin{aufgabe}{Antiproportional – Fehler finden, Begründen, Darstellungswechsel, Anwendung}
\swz[3]{4 Maler brauchen 5 Tage. Kim rechnet: „8 Maler brauchen 10 Tage, denn es sind doppelt so viele.“ Finde den Fehler und rechne richtig.}{}
\swfrage{Warum bleibt bei der Zuordnung Anzahl Arbeiter → Anzahl Tage das Produkt gleich? Begründe.}
\swa{Trage die Punkte der Tabelle ein und verbinde sie zu einer Kurve.}{\wertetabelle[6,3,2,1]{Anzahl Kinder}{Stück je Kind}{1,2,3,6} \begin{ksys}[xmin=0,xmax=7,ymin=0,ymax=7,xlabel=Anzahl Kinder,ylabel=Stück je Kind]\end{ksys}}
\swz[3]{Für das Schmücken der Turnhalle plant die Klasse mit 5 Helfern 3 Stunden. Die Halle ist aber nur 2 Stunden frei. Wie viele Helfer braucht man mindestens?}{}
\end{aufgabe}

\uebersichtskasten{Antiproportional: Je mehr, desto weniger – doppelte Anzahl, halbe Zeit. Das Produkt bleibt gleich. \\ 4 Pumpen brauchen 9 Stunden \quad | links : 4, rechts · 4 \\ 1 Pumpe braucht 36 Stunden \quad | links · 6, rechts : 6 \\ 6 Pumpen brauchen 6 Stunden \quad Probe: 4 · 9 = 36 = 6 · 6}
